\begin{abstract}
A prediction can tell a robot to act before a safety margin runs out.  But a
fixed threshold can also make it act early.  I wanted to understand when
prediction adds something beyond choosing a more conservative threshold.  I
study this question with a simplified timing model and two kinematic simulation
setups, giving both rules the same tuning budget and testing on separate noise
seeds.

In the timing model, the predictive rule is a current-margin threshold whose
buffer scales with the estimated rate of decline.  Extra lead changes the
recovery outcome only within a band set by recovery duration and deadline
slack; old observations consume that lead one for one.  The minimal predictive
horizon wastes less lead than the minimal safe fixed buffer when recovery
needs a similar amount of time across cases but the corresponding margins
differ, provided the rates can be estimated.

The clearest gain appears in the three-robot shared-zone test.  Both tuned
rules succeed in all 1,620 conflict episodes, and prediction reduces mean
completion time by $2.78$\,s (paired 95\% bootstrap interval
$[2.66,2.90]$\,s).  On clear paths, it reduces unnecessary serializations from
1,080 to 21 out of 1,620 episodes.  The two-robot test exposes a different
behavior: the tuned fixed rule records no collisions in 6,750 conflict episodes
and prediction records two.  Together, the tests show where adapting the margin
helps in these simulations and where estimating its rate becomes the weak point.
\end{abstract}

\noindent\textbf{Keywords:} Runtime assurance, multi-robot systems, predictive
monitoring, safety margins, time to collision, reproducible evaluation.
